% sec10_2.tex — Section 10.2: Heat Equation on an Infinite Domain
\section{Heat Equation on an Infinite Domain}
\label{sec:heat-infinite}

We wish to solve the heat equation on the entire real line:
\begin{equation}
\label{eq:heat-inf-pde}
\pd{u}{t} = k\pdd{u}{x}, \quad -\infty < x < \infty
\end{equation}
subject to the initial condition
\begin{equation}
\label{eq:heat-inf-ic}
u(x,0) = f(x).
\end{equation}
We will assume that $u(x,t) \to 0$ as $x \to \pm\infty$ (we will also need $f(x) \to 0$ as $x \to \pm\infty$).

We attempt to use the method of separation of variables.  We look for product solutions of the form $u(x,t) = \phi(x)h(t)$.  The time-dependent part satisfies
\begin{equation}
\label{eq:heat-inf-time}
\od{h}{t} = -\lambda k h,
\end{equation}
and the spatial part satisfies
\begin{equation}
\label{eq:heat-inf-space}
\odd{\phi}{x} = -\lambda\phi.
\end{equation}
The boundedness conditions as $x \to \pm\infty$ replace the usual homogeneous boundary conditions.  Instead of a finite interval eigenvalue problem with a discrete spectrum, we obtain a continuous spectrum.

For \eqref{eq:heat-inf-space}, the condition that $|\phi(x)| < \infty$ as $x \to \pm\infty$ is satisfied for all $\lambda > 0$.  Writing $\lambda = \omega^2$ (with $\omega > 0$), the bounded solutions are $\cos\omega x$ and $\sin\omega x$, or equivalently $e^{\pm i\omega x}$.  Since all positive values of $\lambda$ are allowed, the spectrum is continuous.

The corresponding time-dependent part is
\[
h(t) = e^{-k\omega^2 t}.
\]
Product solutions are $e^{-i\omega x}e^{-k\omega^2 t}$ for all real $\omega$.  A generalized principle of superposition (integrating over the continuous parameter $\omega$ rather than summing over discrete eigenvalues) suggests the solution
\begin{equation}
\label{eq:heat-inf-superposition}
u(x,t) = \int_{-\infty}^{\infty} c(\omega)\,e^{-i\omega x}\,e^{-k\omega^2 t}\,d\omega.
\end{equation}
The initial condition $u(x,0) = f(x)$ is satisfied if
\begin{equation}
\label{eq:heat-inf-ic-integral}
f(x) = \int_{-\infty}^{\infty} c(\omega)\,e^{-i\omega x}\,d\omega.
\end{equation}
This is a Fourier integral representation of $f(x)$; the coefficient $c(\omega)$ is the Fourier transform of $f(x)$.  In the next section we derive the Fourier transform pair and show how $c(\omega)$ can be determined.
